\section{El punto de partida de Self-Attention: Content Addressing y Parallel Reading}

La sección anterior ya contaba con attention, pero solo se usaba en la source--target interface. En self-attention, query, key y value provienen de la misma secuencia, de modo que cada token lee directamente cualquier posición dentro de una sola capa. Al mismo tiempo logra un global receptive field, position-parallel computation y multiplicative interaction; el costo es una score matrix de $n\times n$ y la pérdida de la información de orden.

\subsection{El Attention-like Motif ya existía en otras modalidades}

El curso presenta primero non-local means y texture synthesis para oponerse al relato heroico de que «attention surgió de la nada». El \term{content-based addressing (direccionamiento por contenido)} decide qué posición leer según la similitud entre la query actual y el contenido candidato, no según una dirección fija o un offset; es adecuado para la autosimilitud de las imágenes, la recuperación de memoria y la source alignment.

\lecturefigure{slide-11-self-attention-other-modalities.jpg}{Direccionamiento por contenido en non-local means y texture synthesis}{Diapositiva V011 recuperada del video oficial; 00:17:00}

\begin{knowledgebox}{Lectura de la figura: la vecindad fija y la vecindad por contenido son dos supuestos previos distintos}
A la izquierda, el denoising busca patches similares en toda la imagen; a la derecha, la synthesis busca texturas coincidentes en una base de datos. Primero observe si las posiciones leídas deben ser contiguas y luego quién define la similitud; este mecanismo comparte con attention la estructura de «la query decide la conexión», pero en los algoritmos tradicionales la feature/similarity puede diseñarse a mano, mientras que el Transformer aprende la projection de extremo a extremo.
\end{knowledgebox}

\subsection{Self-Attention: los tokens de una oración se actualizan entre sí}

La subsección anterior explicó por qué el direccionamiento por contenido puede ir más allá de una vecindad fija; aquí se responde cómo se convierte en un cálculo secuencial entrenable de una capa. Al leer el diagrama de estructura que sigue, primero siga cómo un solo token envía una consulta a toda la oración y luego observe por qué las consultas de todos los tokens pueden resolverse en paralelo, por lotes. La \term{self-attention (autoatención)} hace que cada posición de la secuencia produzca query, key y value, y usa el query--key score para ponderar y agregar los value. Los $Q,K,V$ de todas las posiciones pueden generarse con una sola matrix multiplication, y después la score matrix se calcula en paralelo.

\lecturefigure{slide-12-self-attention-structure.jpg}{Self-attention permite que los tokens de una misma secuencia intercambien información directamente}{Diapositiva V012 recuperada del video oficial; 00:17:27}

\begin{knowledgebox}{Lectura de la figura: dentro de una capa el grafo es fully connected, pero los pesos los determinan los datos}
La palabra central puede leer cualquier posición de la oración sin pasar por varios recurrent/convolutional step. Primero observe si cada salida puede conectarse con todas las entradas y luego que distintas query obtienen distintos pesos; fully connected significa que las aristas candidatas existen, no que todas tengan el mismo peso, ni que el modelo ya conozca las relaciones sintácticas.
\end{knowledgebox}

\lecturefigure{slide-13-self-attention-properties.jpg}{Las tres propiedades de diseño de self-attention}{Diapositiva V013 recuperada del video oficial; 00:17:36}

\begin{importantbox}{Tres beneficios y un costo implícito}
Self-attention es paralelizable, captura directamente la token--token interaction y, mediante el dot product, ofrece un multiplicative gating explícito. El costo implícito es que cada par de tokens produce un score, por lo que en la implementación estándar el cálculo y la logit memory crecen con $n^2$; la «conexión total» es a la vez la fuente de su capacidad y el cuello de botella del long-context.
\end{importantbox}

\subsection{El objetivo original también incluía Parallel Writing}

La historia simplificada del Transformer suele decir: «las RNN eran demasiado lentas, así que se cambiaron por attention». El curso añade un objetivo original más ambicioso: se buscaba leer la source en paralelo y también escribir el target en paralelo. La \term{autoregressive decoding (decodificación autorregresiva)} reduce el espacio de salida token por token según $p(y_t\mid y_{<t},x)$; si se predicen todos los tokens a la vez, hay que aprender al mismo tiempo qué tokens pueden ser condicionalmente independientes y qué modes decidir primero.

\lecturefigure{slide-14-transformer-original-motivation.jpg}{Motivación original: leer en paralelo e intentar generar en paralelo}{Diapositiva V014 recuperada del video oficial; 00:19:42}

Las factorization autorregresiva y completamente no autorregresiva son, respectivamente:
\[
p_{\mathrm{AR}}(y\mid x)=\prod_{t=1}^{m}p(y_t\mid y_{<t},x),\qquad
p_{\mathrm{NAR}}(y\mid x)=\prod_{t=1}^{m}p(y_t\mid x).
\]
Aquí $x$ es la entrada, $y$ es la secuencia de salida y $m$ es la longitud de la salida. La segunda fórmula supone que, dado $x$, cada $y_t$ es condicionalmente independiente, lo que dificulta representar las múltiples formulaciones válidas de una traducción y las restricciones entre ellas; en la práctica, los métodos NAR suelen incorporar una latent variable, iterative refinement o un orden por grupos.

\teachervoice{El ponente admite con franqueza que el planteamiento original de parallel decoding no funcionó: en la decodificación autorregresiva, cada elección de palabra «dobla» la distribución de probabilidad y va fijando un mode paso a paso; el modelo no autorregresivo, en cambio, también tiene que aprender a decidir el orden. Al final, el equipo conservó la parallel reading, pero no logró eliminar la sequential generation.}

\begin{warningbox}{Las rutas fallidas también forman parte de la historia de la arquitectura}
El éxito del encoder parallelism no debe reescribirse como si hubiera sido el único objetivo inicial. Registrar el fracaso de la parallel writing explica por qué el decoder todavía necesita una causal mask y por qué la speculative decoding tiene valor práctico, y recuerda a los investigadores que deben distinguir entre el paralelismo teórico y la estructura de dependencias de la salida.
\end{warningbox}

\subsection{Resumen de la sección}

Self-attention convierte la content retrieval, común en muchas modalidades, en el sequence backbone, y obtiene global interaction en una sola capa y aprendizaje de representaciones en paralelo. Al mismo tiempo, la fully parallel generation sigue siendo difícil por el mode ordering y la conditional dependence; el triunfo del Transformer fue un éxito tras un proceso de selección, no la consecución simultánea de todos los objetivos originales.
